\begin{blocksection}
    \question Write a function \lstinline$count_t$, which takes in a dictionary \lstinline$d$ and a string \lstinline$word$.
    The function should count the instances of the letter "t" in \lstinline$word$ and add a key-value pair to the dictionary.
    The key will be \lstinline$word$ and the value will be the number of "t"s in \lstinline$word$
    
    \begin{lstlisting}
    def count_t(d: dict[str, int], word: str) -> None:
        """
        >>> words = {}
        >>> count_t(words, "tatter")
        >>> words["tatter"]
        3
        >>> count_t(words, "tree")
        >>> words
        {'tatter': 3, 'tree': 1}
        """
        _______________________________
        
        for ____________________________:
    
            if ____________________________:
    
                __________________________________
    
        _______________________________
    
    \end{lstlisting}
    
    \begin{solution}
    \begin{lstlisting}
    count = 0
    for c in word:
        if c == 't':
            count += 1
    d[word] = count
    \end{lstlisting}
    \end{solution}
    
\end{blocksection}

\begin{questionmeta}
\begin{blocksection}
\textbf{Teaching Tips}

\textbf{Suggested Time:} 6 min \qquad \textbf{Difficulty:} Easy
\begin{itemize}
    \item \textbf{Read the type hints.}
    \lstinline{dict[str, int]} means string keys and integer values;
    \lstinline{word: str} means a string input. Annotations describe expected
    types; Python does not automatically enforce them.

    \item \textbf{Review dictionary access.}
    A dictionary maps keys to values. Use \lstinline{d[word]} to read a value
    and \lstinline{d[word] = count} to add or update the entry.

    \item \textbf{Plan the loop.}
    Traverse the string directly with \lstinline{for c in word}.
    Start the count at zero, increment it when \lstinline{c == 't'},
    and update the dictionary after the loop.

    \item \textbf{Distinguish mutation from returning.}
    The function changes \lstinline{d} in place and implicitly returns
    \lstinline{None}, matching \lstinline{-> None}.
    Ask students why the doctests inspect \lstinline{words} after the call.
\end{itemize}
\end{blocksection}
\end{questionmeta}
